\documentclass[11pt,a4paper]{article}

% ===== PACKAGES =====
\usepackage[utf8]{inputenc}
\usepackage[T1]{fontenc}
\usepackage[margin=1in]{geometry}
\usepackage{hyperref}
\usepackage{listings}
\usepackage{xcolor}
\usepackage{booktabs}
\usepackage{array}
\usepackage{longtable}
\usepackage{fancyhdr}
\usepackage{titlesec}
\usepackage{tcolorbox}
\usepackage{fontawesome5}
\usepackage{amssymb}
\usepackage{amsmath}
\usepackage{enumitem}

% ===== COLORS =====
\definecolor{codegreen}{rgb}{0.25,0.49,0.48}
\definecolor{codegray}{rgb}{0.5,0.5,0.5}
\definecolor{codepurple}{rgb}{0.58,0,0.82}
\definecolor{codeblue}{rgb}{0.13,0.29,0.53}
\definecolor{codeorange}{rgb}{0.8,0.4,0.0}
\definecolor{backcolour}{rgb}{0.95,0.95,0.95}
\definecolor{warningbg}{rgb}{1.0,0.95,0.85}
\definecolor{warningborder}{rgb}{0.9,0.6,0.2}
\definecolor{tipbg}{rgb}{0.9,0.97,0.9}
\definecolor{tipborder}{rgb}{0.3,0.7,0.3}
\definecolor{notebg}{rgb}{0.9,0.93,1.0}
\definecolor{noteborder}{rgb}{0.3,0.4,0.7}

% ===== CODE LISTING STYLE =====
\lstdefinelanguage{IEC61131}{
    keywords={VAR, END_VAR, VAR_INPUT, VAR_OUTPUT, VAR_IN_OUT, VAR_GLOBAL, 
              TYPE, END_TYPE, STRUCT, END_STRUCT, FUNCTION, END_FUNCTION,
              FUNCTION_BLOCK, END_FUNCTION_BLOCK, PROGRAM, END_PROGRAM,
              IF, THEN, ELSE, ELSIF, END_IF, CASE, OF, END_CASE,
              FOR, TO, BY, DO, END_FOR, WHILE, END_WHILE, REPEAT, UNTIL, END_REPEAT,
              TRUE, FALSE, AND, OR, NOT, XOR, MOD, RETURN,
              ARRAY, OF, POINTER, REFERENCE, TO, AT, EXTENDS, IMPLEMENTS,
              METHOD, END_METHOD, PROPERTY, END_PROPERTY,
              INTERFACE, END_INTERFACE,
              PUBLIC, PRIVATE, PROTECTED, INTERNAL, FINAL, ABSTRACT,
              THIS, SUPER, VAR_STAT, CONSTANT, VAR_INST},
    keywordstyle=\color{codeblue}\bfseries,
    keywords=[2]{BOOL, BYTE, WORD, DWORD, LWORD, SINT, USINT, INT, UINT, 
                 DINT, UDINT, LINT, ULINT, REAL, LREAL, STRING, WSTRING,
                 TIME, LTIME, DATE, TIME_OF_DAY, TOD, DATE_AND_TIME, DT,
                 TcEventSeverity, T_FILETIME, E_Color},
    keywordstyle=[2]\color{codepurple}\bfseries,
    keywords=[3]{ADR, SIZEOF, REF, FILETIME_TO_DT, GETSYSTEMTIME, FB_init,
                 CONCAT, TO_STRING, FB_FileOpen, FB_FilePuts, FB_FileClose},
    keywordstyle=[3]\color{codeorange}\bfseries,
    sensitive=false,
    morecomment=[l]{//},
    morecomment=[s]{(*}{*)},
    commentstyle=\color{codegreen}\itshape,
    stringstyle=\color{codeorange},
    morestring=[b]',
    morestring=[b]",
}

\lstdefinestyle{iecstyle}{
    language=IEC61131,
    backgroundcolor=\color{backcolour},
    basicstyle=\ttfamily\small,
    breakatwhitespace=false,
    breaklines=true,
    captionpos=b,
    keepspaces=true,
    numbers=left,
    numbersep=8pt,
    numberstyle=\tiny\color{codegray},
    showspaces=false,
    showstringspaces=false,
    showtabs=false,
    tabsize=4,
    frame=single,
    framerule=0.5pt,
    rulecolor=\color{codegray},
    xleftmargin=15pt,
    framexleftmargin=15pt,
}

\lstset{style=iecstyle}

% ===== TCOLORBOX STYLES =====
\tcbuselibrary{skins,breakable}

\newtcolorbox{warningbox}{
    colback=warningbg,
    colframe=warningborder,
    boxrule=1pt,
    left=10pt,
    right=10pt,
    top=8pt,
    bottom=8pt,
    breakable,
    before upper={\faExclamationTriangle\ \textbf{Warning:} }
}

\newtcolorbox{tipbox}{
    colback=tipbg,
    colframe=tipborder,
    boxrule=1pt,
    left=10pt,
    right=10pt,
    top=8pt,
    bottom=8pt,
    breakable,
    before upper={\faLightbulb\ \textbf{Tip:} }
}

\newtcolorbox{notebox}{
    colback=notebg,
    colframe=noteborder,
    boxrule=1pt,
    left=10pt,
    right=10pt,
    top=8pt,
    bottom=8pt,
    breakable,
    before upper={\faInfoCircle\ \textbf{Note:} }
}

% ===== HEADER/FOOTER =====
\pagestyle{fancy}
\fancyhf{}
\fancyhead[L]{\leftmark}
\fancyhead[R]{TwinCAT 3 Lecture Notes}
\fancyfoot[C]{\thepage}
\renewcommand{\headrulewidth}{0.4pt}
\renewcommand{\footrulewidth}{0.4pt}

% ===== HYPERREF SETUP =====
\hypersetup{
    colorlinks=true,
    linkcolor=codeblue,
    filecolor=magenta,
    urlcolor=cyan,
    pdftitle={TwinCAT 3 - Program Flow Instructions and State Machines},
    pdfauthor={},
    bookmarks=true,
}

% ===== TITLE FORMATTING =====
\titleformat{\section}
    {\normalfont\Large\bfseries}{\thesection}{1em}{}[{\titlerule[0.8pt]}]
\titleformat{\subsection}
    {\normalfont\large\bfseries}{\thesubsection}{1em}{}
\titleformat{\subsubsection}
    {\normalfont\normalsize\bfseries}{\thesubsubsection}{1em}{}

% ===== DOCUMENT INFO =====
\title{
    \vspace{-1cm}
    \textbf{Lecture Notes: TwinCAT 3}\\
    \Large Program Flow Instructions \& State Machines
}
\author{}
\date{\today}

% ===== BEGIN DOCUMENT =====
\begin{document}

\maketitle

\tableofcontents
\newpage

% ============================================================
\section{Introduction to Program Flow Control}
% ============================================================

Instructions are used to control the flow of execution within a PLC program. There are four primary types of instructions covered:

\begin{itemize}[leftmargin=*]
    \item \textbf{Conditional Statements:} IF/ELSE and CASE switches.
    \item \textbf{Iteration Statements (Loops):} FOR loops and WHILE loops.
\end{itemize}

\begin{table}[h]
\centering
\begin{tabular}{@{}lll@{}}
\toprule
\textbf{Category} & \textbf{Instruction} & \textbf{Use Case} \\
\midrule
Conditional & IF/ELSIF/ELSE & Boolean logic decisions \\
Conditional & CASE & Multiple discrete value testing \\
Iteration & FOR & Fixed number of iterations \\
Iteration & WHILE & Condition-based iteration \\
\bottomrule
\end{tabular}
\caption{Program Flow Instructions Overview}
\end{table}

% ============================================================
\section{Conditional Statements}
% ============================================================

\subsection{IF / ELSIF / ELSE}

\begin{itemize}[leftmargin=*]
    \item \textbf{IF:} Specifies a block of code to execute only if a boolean expression evaluates to \textbf{TRUE}.
    
    \item \textbf{ELSE:} Provides an alternative path of execution if the IF condition is \textbf{FALSE}.
    
    \item \textbf{ELSIF:} Allows for testing additional conditions if the preceding conditions are false. You can use as many ELSIFs as needed.
\end{itemize}

\subsubsection*{Example: IF/ELSE Logic}

\begin{lstlisting}
// Program: Temperature Classification
// Demonstrates IF/ELSIF/ELSE conditional logic
PROGRAM MAIN
VAR
    nTemperature    : INT;           // Current temperature reading
    sStatus         : STRING(50);    // Status message
    bAlarmActive    : BOOL;          // Alarm flag
    
    // Threshold constants
    nHighThreshold  : INT := 80;     // High temperature threshold
    nLowThreshold   : INT := 20;     // Low temperature threshold
END_VAR

// IF/ELSIF/ELSE structure for temperature classification
IF nTemperature > nHighThreshold THEN
    // Temperature exceeds high threshold - critical condition
    sStatus := 'CRITICAL: Temperature too high!';
    bAlarmActive := TRUE;
    
ELSIF nTemperature < nLowThreshold THEN
    // Temperature below low threshold - warning condition
    sStatus := 'WARNING: Temperature too low!';
    bAlarmActive := TRUE;
    
ELSIF nTemperature = 50 THEN
    // Temperature exactly at optimal value
    sStatus := 'OPTIMAL: Perfect temperature';
    bAlarmActive := FALSE;
    
ELSE
    // Temperature within normal operating range
    sStatus := 'NORMAL: Temperature acceptable';
    bAlarmActive := FALSE;
    
END_IF
\end{lstlisting}

\begin{lstlisting}
// Additional example: Nested IF statements
VAR
    bSensorOK       : BOOL;
    bSystemReady    : BOOL;
    nErrorCode      : INT;
END_VAR

IF bSensorOK THEN
    IF bSystemReady THEN
        // Both conditions met - proceed with operation
        nErrorCode := 0;
    ELSE
        // Sensor OK but system not ready
        nErrorCode := 2;
    END_IF
ELSE
    // Sensor failure - critical error
    nErrorCode := 1;
END_IF
\end{lstlisting}

\subsection{CASE Statements}

\begin{itemize}[leftmargin=*]
    \item \textbf{Purpose:} Used to perform different actions based on the specific value of a condition variable (typically an \textbf{INT} or \textbf{WORD}).
    
    \item \textbf{Note:} Strings are \textbf{not} supported in CASE statements; for strings, you must use IF/ELSIF.
    
    \item \textbf{Preferred Use:} Generally preferred over long chains of IF/ELSIF when testing many different values of a single variable.
\end{itemize}

\subsubsection*{Example: CASE Statement}

\begin{lstlisting}
// Program: State-Based Control using CASE
// Demonstrates CASE statement with integer state variable
PROGRAM MAIN
VAR
    nState          : INT := 1;      // Current state (1, 2, 3, etc.)
    sStateDesc      : STRING(50);    // State description
    bMotorEnable    : BOOL;          // Motor control output
    bValveOpen      : BOOL;          // Valve control output
    fSetpoint       : REAL;          // Process setpoint
END_VAR

// CASE statement for state machine control
CASE nState OF

    1:  // Initialization State
        sStateDesc := 'State 1: Initialization';
        bMotorEnable := FALSE;
        bValveOpen := FALSE;
        fSetpoint := 0.0;
        
        // Transition logic would go here
        // IF bInitComplete THEN nState := 2; END_IF
        
    2:  // Startup State
        sStateDesc := 'State 2: System Startup';
        bMotorEnable := TRUE;
        bValveOpen := FALSE;
        fSetpoint := 25.0;
        
    3:  // Running State
        sStateDesc := 'State 3: Normal Operation';
        bMotorEnable := TRUE;
        bValveOpen := TRUE;
        fSetpoint := 75.0;
        
    4, 5:  // Multiple values can share the same action
        sStateDesc := 'State 4 or 5: Maintenance Mode';
        bMotorEnable := FALSE;
        bValveOpen := FALSE;
        
    10..20:  // Range of values (10 through 20)
        sStateDesc := 'States 10-20: Error Recovery';
        bMotorEnable := FALSE;
        bValveOpen := FALSE;
        fSetpoint := 0.0;

ELSE
    // Default case - handles any undefined state values
    sStateDesc := 'UNKNOWN STATE - Check nState value';
    bMotorEnable := FALSE;
    bValveOpen := FALSE;
    fSetpoint := 0.0;
    
END_CASE
\end{lstlisting}

\begin{warningbox}
CASE statements do not support STRING variables. If you need to branch based on string values, you must use IF/ELSIF chains or convert the string to an enumeration or integer first.
\end{warningbox}

% ============================================================
\section{Iteration Statements (Loops)}
% ============================================================

\subsection{FOR Loops}

\begin{itemize}[leftmargin=*]
    \item \textbf{Function:} Executes a block of code a specific number of times.
    
    \item \textbf{Use Case:} Highly efficient for iterating through \textbf{arrays}.
    
    \item \textbf{Mechanism:} Uses a counter variable that automatically increments by 1 (or a custom step size defined by the \texttt{BY} keyword) at the end of each loop until a target value is reached.
\end{itemize}

\subsubsection*{Example: FOR Loop Array Summation}

\begin{lstlisting}
// Program: Array Processing with FOR Loop
// Demonstrates summing array elements using a FOR loop
PROGRAM MAIN
VAR
    // Array of data values (indices 1 to 5)
    aData   : ARRAY[1..5] OF INT := [10, 20, 30, 40, 50];
    
    // Loop counter variable
    i       : INT;
    
    // Accumulator for sum
    nSum    : INT := 0;
    
    // Additional example variables
    nMax    : INT;
    nMin    : INT;
    fAverage: REAL;
END_VAR

// Reset sum before calculation
nSum := 0;

// FOR loop to sum all elements in the array
// Loop runs from index 1 to index 5 (inclusive)
FOR i := 1 TO 5 DO
    nSum := nSum + aData[i];
END_FOR

// Result: nSum = 10 + 20 + 30 + 40 + 50 = 150
\end{lstlisting}

\begin{lstlisting}
// Additional FOR loop examples

// Example 1: Find maximum value in array
nMax := aData[1];  // Initialize with first element
FOR i := 2 TO 5 DO
    IF aData[i] > nMax THEN
        nMax := aData[i];
    END_IF
END_FOR

// Example 2: FOR loop with custom step size (BY keyword)
// Count down from 10 to 0 by steps of 2
VAR
    nCountdown : INT;
    j          : INT;
END_VAR

FOR j := 10 TO 0 BY -2 DO
    nCountdown := j;  // Values: 10, 8, 6, 4, 2, 0
END_FOR

// Example 3: Initialize array with calculated values
VAR
    aSquares : ARRAY[0..9] OF INT;
    k        : INT;
END_VAR

FOR k := 0 TO 9 DO
    aSquares[k] := k * k;  // Store square of index
END_FOR
// Result: aSquares = [0, 1, 4, 9, 16, 25, 36, 49, 64, 81]
\end{lstlisting}

\subsection{WHILE Loops}

\begin{itemize}[leftmargin=*]
    \item \textbf{Function:} Continues to execute a block of code as long as a specified condition remains \textbf{TRUE}.
    
    \item \textbf{Risk:} Unlike FOR loops, the developer must ensure the condition eventually becomes false to avoid an infinite loop that could crash the PLC cycle.
\end{itemize}

\begin{lstlisting}
// Example: WHILE loop for searching an array
VAR
    aBuffer     : ARRAY[0..99] OF INT;
    nSearchVal  : INT := 42;
    nFoundIndex : INT := -1;    // -1 indicates not found
    nIndex      : INT := 0;
    bFound      : BOOL := FALSE;
END_VAR

// Search for value using WHILE loop
nIndex := 0;
bFound := FALSE;

WHILE (nIndex <= 99) AND (NOT bFound) DO
    IF aBuffer[nIndex] = nSearchVal THEN
        nFoundIndex := nIndex;
        bFound := TRUE;
    ELSE
        nIndex := nIndex + 1;  // CRITICAL: Must increment to avoid infinite loop!
    END_IF
END_WHILE
\end{lstlisting}

\begin{warningbox}
WHILE loops can cause infinite loops if the exit condition is never met. Always ensure that the loop condition will eventually become FALSE, and consider adding a maximum iteration counter as a safety mechanism.
\end{warningbox}

% ============================================================
\section{State Machines and Best Practices}
% ============================================================

In PLC programming, \textbf{state machines} are a standard way to manage complex sequences.

\begin{itemize}[leftmargin=*]
    \item \textbf{Avoiding ``Magic Numbers'':} Many programmers use integers (10, 20, 30) for states. The sources strongly advise \textbf{against} this because numbers do not describe what the code is doing.
    
    \item \textbf{Local Enumerations:} Use a \textbf{local ENUM} (defined within the function block) to name states (e.g., \texttt{FILE\_OPEN}, \texttt{WAIT\_FOR\_EVENT}). This makes the code self-descriptive and prevents ``namespace pollution''.
\end{itemize}

\begin{table}[h]
\centering
\begin{tabular}{@{}lll@{}}
\toprule
\textbf{Approach} & \textbf{Example} & \textbf{Readability} \\
\midrule
Magic Numbers & \texttt{CASE nState OF 10: ...} & Poor \\
Global ENUM & \texttt{CASE eState OF E\_State.Running: ...} & Good \\
Local ENUM & \texttt{CASE eWriteState OF E\_WriteState.FILE\_OPEN: ...} & Best \\
\bottomrule
\end{tabular}
\caption{State Machine Implementation Approaches}
\end{table}

\subsubsection*{Example: Local ENUM State Machine}

\begin{lstlisting}
// Function Block: FB_CsvEventWriter
// Demonstrates local ENUM for state machine implementation
FUNCTION_BLOCK FB_CsvEventWriter
VAR
    // LOCAL ENUMERATION - defined within the function block
    // Prevents namespace pollution and makes code self-descriptive
    {attribute 'strict'}
    E_WriteState : (
        IDLE,                   // Waiting for write request
        FILE_OPEN_TRIGGER,      // Trigger file open operation
        FILE_OPEN_WAIT,         // Wait for file to open
        WRITE_DATA_TRIGGER,     // Trigger data write
        WRITE_DATA_WAIT,        // Wait for write completion
        FILE_CLOSE_TRIGGER,     // Trigger file close
        FILE_CLOSE_WAIT,        // Wait for file to close
        WAIT_FOR_EVENT,         // Wait for next event to log
        DONE,                   // Operation completed successfully
        ERROR                   // Error state
    );
    
    // Current state variable using the local enum
    eWriteState : E_WriteState := E_WriteState.IDLE;
    
    // File operation function blocks
    fbFileOpen  : FB_FileOpen;
    fbFilePuts  : FB_FilePuts;
    fbFileClose : FB_FileClose;
    
    // File handle and path
    hFile       : UINT;
    sFilePath   : STRING(255) := 'C:\EventLogs\events.csv';
    
    // Status flags
    bBusy       : BOOL;
    bError      : BOOL;
    nErrorId    : UDINT;
END_VAR
\end{lstlisting}

\begin{lstlisting}
// State machine implementation using CASE with local ENUM
CASE eWriteState OF

    E_WriteState.IDLE:
        // Initial state - waiting for trigger
        bBusy := FALSE;
        bError := FALSE;
        // Transition: IF bStartWrite THEN eWriteState := E_WriteState.FILE_OPEN_TRIGGER; END_IF
        
    E_WriteState.FILE_OPEN_TRIGGER:
        // Trigger the file open operation
        bBusy := TRUE;
        fbFileOpen(
            bExecute := TRUE,
            sPathName := sFilePath,
            nMode := FOPEN_MODEAPPEND OR FOPEN_MODETEXT
        );
        eWriteState := E_WriteState.FILE_OPEN_WAIT;
        
    E_WriteState.FILE_OPEN_WAIT:
        // Wait for file open to complete
        fbFileOpen(bExecute := FALSE);
        
        IF NOT fbFileOpen.bBusy THEN
            IF fbFileOpen.bError THEN
                nErrorId := fbFileOpen.nErrId;
                eWriteState := E_WriteState.ERROR;
            ELSE
                hFile := fbFileOpen.hFile;
                eWriteState := E_WriteState.WRITE_DATA_TRIGGER;
            END_IF
        END_IF
        
    E_WriteState.WRITE_DATA_TRIGGER:
        // Trigger data write operation
        fbFilePuts(
            bExecute := TRUE,
            hFile := hFile,
            sLine := sCsvLine  // CSV formatted string
        );
        eWriteState := E_WriteState.WRITE_DATA_WAIT;
        
    E_WriteState.WRITE_DATA_WAIT:
        // Wait for write to complete
        fbFilePuts(bExecute := FALSE);
        
        IF NOT fbFilePuts.bBusy THEN
            IF fbFilePuts.bError THEN
                nErrorId := fbFilePuts.nErrId;
                eWriteState := E_WriteState.ERROR;
            ELSE
                eWriteState := E_WriteState.FILE_CLOSE_TRIGGER;
            END_IF
        END_IF
        
    E_WriteState.FILE_CLOSE_TRIGGER:
        // Trigger file close to flush data to disk
        fbFileClose(bExecute := TRUE, hFile := hFile);
        eWriteState := E_WriteState.FILE_CLOSE_WAIT;
        
    E_WriteState.FILE_CLOSE_WAIT:
        // Wait for file close to complete
        fbFileClose(bExecute := FALSE);
        
        IF NOT fbFileClose.bBusy THEN
            IF fbFileClose.bError THEN
                nErrorId := fbFileClose.nErrId;
                eWriteState := E_WriteState.ERROR;
            ELSE
                eWriteState := E_WriteState.DONE;
            END_IF
        END_IF
        
    E_WriteState.DONE:
        // Operation completed successfully
        bBusy := FALSE;
        eWriteState := E_WriteState.IDLE;
        
    E_WriteState.ERROR:
        // Error handling state
        bBusy := FALSE;
        bError := TRUE;
        // Remain in error state until acknowledged
        
END_CASE
\end{lstlisting}

\begin{tipbox}
Using the \texttt{\{attribute 'strict'\}} pragma on local enumerations ensures type safety and prevents accidental assignment of integer values directly to the enum variable.
\end{tipbox}

% ============================================================
\section{Practical Implementation: CSV Persistent Storage}
% ============================================================

The lecture demonstrates these concepts by building a CSV event logger using a state machine and built-in Beckhoff function blocks (\texttt{FB\_FileOpen}, \texttt{FB\_FilePuts}, \texttt{FB\_FileClose}).

\subsection{Asynchronous Communication}

\begin{itemize}[leftmargin=*]
    \item \textbf{Windows Dependency:} File operations are handled by Windows, not the real-time TwinCAT kernel. This communication is \textbf{asynchronous} and non-deterministic.
    
    \item \textbf{Trigger and Wait Pattern:} Because an operation might take several PLC cycles to complete, states must be split:
    \begin{enumerate}
        \item \textbf{Trigger Step:} Set \texttt{bExecute := TRUE} to start the operation.
        \item \textbf{Wait Step:} Continuously call the FB and check the \texttt{bBusy} and \texttt{bError} flags until the operation is complete.
    \end{enumerate}
\end{itemize}

\begin{figure}[h]
\centering
\begin{tabular}{|c|c|c|}
\hline
\textbf{Step} & \textbf{Action} & \textbf{Next Step Condition} \\
\hline
Trigger & Set \texttt{bExecute := TRUE} & Immediate \\
\hline
Wait & Check \texttt{bBusy} flag & \texttt{bBusy = FALSE} \\
\hline
Evaluate & Check \texttt{bError} flag & Proceed or Error state \\
\hline
\end{tabular}
\caption{Asynchronous Operation Pattern}
\end{figure}

\subsection{Helper Methods and String Manipulation}

\begin{itemize}[leftmargin=*]
    \item \textbf{CreateCsvString():} A private helper method used to convert a data structure into a comma-separated string.
    
    \item \textbf{TO\_STRING Pragma:} Adding \texttt{\{attribute 'to\_string'\}} to an enumeration allows TwinCAT to automatically convert the ENUM name into a string, which is useful for logging.
    
    \item \textbf{CONCAT:} A function used to join multiple strings together with commas and special characters like \texttt{\$n} (new line).
\end{itemize}

\subsubsection*{Example: CSV String Building}

\begin{lstlisting}
// Enumeration with 'to_string' attribute for automatic string conversion
{attribute 'qualified_only'}
{attribute 'to_string'}
TYPE E_EventType :
(
    Info := 0,
    Warning := 1,
    Error := 2,
    Critical := 3
);
END_TYPE

{attribute 'qualified_only'}
{attribute 'to_string'}
TYPE E_EventSeverity :
(
    Low := 0,
    Medium := 1,
    High := 2
);
END_TYPE
\end{lstlisting}

\begin{lstlisting}
// Private Method: CreateCsvString
// Converts an event structure into a comma-separated string for file storage
METHOD PRIVATE CreateCsvString : STRING(1000)
VAR_INPUT
    stEvent : REFERENCE TO ST_Event;  // Event to convert
END_VAR
VAR
    sTempString   : STRING(1000);     // Temporary string builder
    sEventType    : STRING(50);       // Event type as string
    sSeverity     : STRING(50);       // Severity as string
    sTimestamp    : STRING(50);       // Timestamp as string
    sEventId      : STRING(20);       // Event ID as string
END_VAR

// Convert enumeration values to strings using TO_STRING
// (Requires {attribute 'to_string'} on the enum definition)
sEventType := TO_STRING(stEvent.eEventType);
sSeverity := TO_STRING(stEvent.eEventSeverity);

// Convert numeric ID to string
sEventId := UDINT_TO_STRING(stEvent.nIdentity);

// Convert timestamp to string
sTimestamp := DT_TO_STRING(stEvent.dtTimestamp);

// Build CSV string using CONCAT
// Format: EventType,Severity,EventID,Timestamp,Description
sTempString := '';

// Add event type
sTempString := CONCAT(sTempString, sEventType);

// Add comma separator and severity
sTempString := CONCAT(sTempString, ',');
sTempString := CONCAT(sTempString, sSeverity);

// Add comma separator and event ID
sTempString := CONCAT(sTempString, ',');
sTempString := CONCAT(sTempString, sEventId);

// Add comma separator and timestamp
sTempString := CONCAT(sTempString, ',');
sTempString := CONCAT(sTempString, sTimestamp);

// Add comma separator and description text
sTempString := CONCAT(sTempString, ',');
sTempString := CONCAT(sTempString, stEvent.sText);

// Add newline character for end of CSV row
// $n is the escape sequence for newline in IEC 61131-3
sTempString := CONCAT(sTempString, '$n');

// Return the completed CSV string
CreateCsvString := sTempString;
\end{lstlisting}

\begin{lstlisting}
// Example output from CreateCsvString:
// "Warning,High,1001,DT#2024-06-15-14:30:45,Pump pressure exceeded threshold$n"

// Alternative: More compact CONCAT chaining
sTempString := CONCAT(
    CONCAT(
        CONCAT(
            CONCAT(
                CONCAT(
                    CONCAT(sEventType, ','),
                    sSeverity
                ),
                ','
            ),
            sEventId
        ),
        ','
    ),
    CONCAT(
        CONCAT(sTimestamp, ','),
        CONCAT(stEvent.sText, '$n')
    )
);
\end{lstlisting}

\subsection{Important Discovery}

\begin{itemize}[leftmargin=*]
    \item \textbf{File Buffering:} In TwinCAT, data may not be written to the file immediately. The implementation found that explicitly calling \texttt{FB\_FileClose} after writing an event was necessary to ensure the data was flushed to the disk.
\end{itemize}

\begin{warningbox}
File write operations in TwinCAT are buffered by Windows. To ensure data is actually written to disk (especially important for logging critical events), you must call \texttt{FB\_FileClose} after writing. Simply calling \texttt{FB\_FilePuts} does not guarantee the data reaches the physical file.
\end{warningbox}

\begin{lstlisting}
// Critical pattern: Always close file to flush buffer
// This ensures data is written to disk

// After write operation completes...
E_WriteState.FILE_CLOSE_TRIGGER:
    // IMPORTANT: Close file to flush write buffer to disk
    // Without this, data may remain in memory buffer and not be saved
    fbFileClose(bExecute := TRUE, hFile := hFile);
    eWriteState := E_WriteState.FILE_CLOSE_WAIT;
\end{lstlisting}

% ============================================================
\section*{Quick Reference Card}
\addcontentsline{toc}{section}{Quick Reference Card}
% ============================================================

\subsection*{Conditional Statement Syntax}

\begin{lstlisting}
// IF/ELSIF/ELSE
IF <condition1> THEN
    // Code block 1
ELSIF <condition2> THEN
    // Code block 2
ELSE
    // Default code block
END_IF

// CASE Statement
CASE <variable> OF
    <value1>:
        // Code for value1
    <value2>, <value3>:
        // Code for value2 or value3
    <start>..<end>:
        // Code for range
ELSE
    // Default code
END_CASE
\end{lstlisting}

\subsection*{Loop Syntax}

\begin{lstlisting}
// FOR Loop
FOR <counter> := <start> TO <end> BY <step> DO
    // Loop body
END_FOR

// WHILE Loop
WHILE <condition> DO
    // Loop body (must modify condition!)
END_WHILE
\end{lstlisting}

\subsection*{State Machine Best Practices}

\begin{table}[h]
\centering
\begin{tabular}{@{}ll@{}}
\toprule
\textbf{Practice} & \textbf{Recommendation} \\
\midrule
State Values & Use local ENUMs, not magic numbers \\
State Names & Descriptive (e.g., \texttt{FILE\_OPEN\_WAIT}) \\
Async Operations & Split into TRIGGER and WAIT states \\
Error Handling & Include dedicated ERROR state \\
File Operations & Always close files to flush buffers \\
\bottomrule
\end{tabular}
\caption{State Machine Best Practices}
\end{table}

\subsection*{String Functions}

\begin{table}[h]
\centering
\begin{tabular}{@{}ll@{}}
\toprule
\textbf{Function} & \textbf{Purpose} \\
\midrule
\texttt{CONCAT(s1, s2)} & Join two strings \\
\texttt{TO\_STRING(enum)} & Convert ENUM to string (requires attribute) \\
\texttt{INT\_TO\_STRING(n)} & Convert integer to string \\
\texttt{DT\_TO\_STRING(dt)} & Convert DATE\_AND\_TIME to string \\
\texttt{\$n} & Newline escape sequence \\
\bottomrule
\end{tabular}
\caption{Common String Functions}
\end{table}

\vfill

\begin{center}
\textit{Last Updated: \today}
\end{center}

\end{document}